% Generated draft: requires mathematical and source review.
\documentclass[11pt]{article}
\usepackage[margin=25mm]{geometry}
\usepackage{fontspec}
\setmainfont{Latin Modern Roman}
\usepackage{amsmath,amssymb,mathtools,unicode-math}
\setmathfont{Latin Modern Math}
\usepackage{enumitem,needspace,xurl,hyperref}
\setlength{\parindent}{0pt}
\setlength{\parskip}{5pt}
\setlength{\emergencystretch}{3em}
\newcommand{\N}{\mathbb{N}}
\newcommand{\Z}{\mathbb{Z}}
\newcommand{\Q}{\mathbb{Q}}
\newcommand{\R}{\mathbb{R}}
\newcommand{\C}{\mathbb{C}}
\newcommand{\F}{\mathbb{F}}
\newcommand{\A}{\mathbb{A}}
\newcommand{\PP}{\mathbb{P}}
\newcommand{\E}{\mathbb{E}}
\newcommand{\eps}{\varepsilon}
\newcommand{\dd}{\,\mathrm{d}}
\begin{document}
\section*{UnsolvedMath ID: 30005040}
\textbf{Source problem code:} catalogue entry
\subsection{Definitions and mathematical statement}
For the word-metric part of the problem, let \(S\subseteq\Z\) be an infinite generating set of the additive group \(\Z\), so that \(\langle S\rangle=\Z\). No symmetry assumption on \(S\) is necessary: put \(S^{\pm}=S\cup(-S)\) and define
\[
\ell_S(k)=\min\left\{m\in\Z_{\geq0}: k=s_1+\cdots+s_m,\ s_i\in S^{\pm}\right\},\qquad d_S(x,y)=\ell_S(y-x).
\]
The empty sum is \(0\). This specifies the usual symmetric word-metric convention; the supplied excerpt does not explicitly state its convention.

A coarse structure on a set \(X\) is a family of subsets of \(X\times X\), called entourages, containing the diagonal and closed under taking subsets, finite unions, inverses, and compositions of relations. For a metric \(d\), its metric coarse structure is
\[
\mathcal E_d=\left\{E\subseteq X\times X:\text{there exists }R\in[0,\infty)\text{ such that }d(x,y)\leq R\text{ for every }(x,y)\in E\right\}.
\]
Write \(d_0(x,y)=|x-y|\). Under the usual inclusion convention, a coarsification of \((\Z,\mathcal E_{d_0})\) is a coarse structure \(\mathcal E\) on \(\Z\) with \(\mathcal E_{d_0}\subseteq\mathcal E\). The supplied report discusses coarse groups, but does not supply their definitions or any additional compatibility requirements. Consequently this convention makes the metric subquestions precise, while the full intended class of coarsifications remains qualified by that missing context.

The requested problem is to classify the coarsifications of \(\Z\). In particular:
\begin{enumerate}
\item For every infinite generating set \(S\subseteq\Z\), determine whether \(d_S\) is bounded, meaning
\[
\exists D\in[0,\infty)\ \forall x,y\in\Z,\qquad d_S(x,y)\leq D.
\]
\item For every pair of infinite generating sets \(S,T\subseteq\Z\), determine whether they induce the same coarse structure, interpreted literally on the underlying set \(\Z\) as
\[
\mathcal E_{d_S}=\mathcal E_{d_T}.
\]
\end{enumerate}
The equivalence relation intended for the overall classification is not specified; equality of coarse structures and equivalence by maps must not be silently identified.

The source also includes coarsifications arising from group topologies \(\tau\) on \(\Z\), written \((\Z,\tau)\), without specifying how a topology determines the relevant coarse structure. It reports that profinite topologies yield at least a continuum, \(2^{\aleph_0}\), of different coarsifications, citing its reference [3]. Its subsequent question about a larger number of distinct coarsifications has corrupted cardinal notation: the extraction contains \texttt{22} and a displaced \(\aleph_0\). A possible reading is \(2^{2^{\aleph_0}}\), but the supplied text does not reliably establish that reading or the meaning of distinctness.

For \(S\) equal to the set of positive primes, with the symmetric word convention above, the source describes boundedness of \(d_S\) as a non-effective version of the Goldbach conjecture. This is retained as a source qualification; no particular formulation of Goldbach or precise equivalence is supplied here. The registry labels the problem open, which is a supplied status rather than an independently verified conclusion.
\subsection{Short English statement}
Classify the coarse structures on \(\Z\) intended by the source. In particular, characterize which infinite generating sets give bounded word metrics and which pairs give exactly the same metric coarse structure. The broader classification requires clarification of the source's coarse-group conventions, its construction from group topologies, its equivalence relation, and a corrupted cardinality question.
\subsection{Sources}
\begin{itemize}
\item Local registry: \texttt{data/unsolvedmath/problems.json}, permanent UnsolvedMath ID \(30005040\), title \emph{Classification of Coarsifications of the Integers}, exact source code \texttt{catalogue entry}. The supplied record identifies the extracted question as Problem 7 and supplies no individual proposer.
\item Supplied report citation: \emph{Geometric Structures in Group Theory (2023)}, Oberwolfach Reports, pp.~517--576, DOI \texttt{10.4171/owr/2022/11}: \url{https://doi.org/10.4171/owr/2022/11}. The record gives proposed year 2022. The relation between that year and the citation's 2023 designation requires source review.
\item Supplied literature source, labelled \emph{Oberwolfach report}: \url{https://doi.org/10.4171/owr/2022/11}.
\item Supplied literature source, labelled \emph{Coarse invertibility spectrum}: \url{https://doi.org/10.1007/s00209-025-03764-0}. No authors or further bibliographic details are supplied.
\item Literature assessment dated exactly \texttt{2026-08-22}: the supplied curation labels the classification problem open and states that recent work distinguishes some examples without providing a complete classification. This is supplied dated curation, not independently verified. The background explicitly calls for independent verification of time-sensitive status and cited claims.
\item The registry labels the statement \texttt{corrected\_verified} and reports that the official report was checked directly. These are supplied provenance claims; no source visit or independent verification is asserted in this draft. The excerpt's references [3] and [5], and [2] in adjacent material, lack bibliographic entries in the supplied record.
\end{itemize}
\end{document}
